\bookStart{The Knosping Galder}\setBookCode{Chnospinci}

\begin{flushright}%
\textbf{Dating:} C9th--10th

\textbf{Meter:} \Fornyrdislag%para
\end{flushright}

\section{Introduction}

The \textbf{Knospink galder} is a short OHG galder for blessing a house.

\subsection{Preservation}

The galder is found in Ms. Car. C 176, presently in Zürich.  It is found on fol.~154r amidst a number of magical charms in Latin.  It is given the title \emph{Ad signandum domum contra diabolum} ‘To bless a house against devils’.

\subsection{Interpretation}

The galder begins by addressing the “wight” explicitly, recognising it as such.  This recognition is the equivalent of the naming of it and brings the dæmonic spirit out into the open, vulnerable to the effects of the exorcism.

The titular \emph{chnospinci} ‘Knosping’ serves to drive away the wight.  It is undoubtedly a magic word and has no clear meaning.  Metrically it appears that the suffix \emph{-pinci} has secondary stress which suggests the segmentation \emph{chnos-pinci}, though this does not help much with its etymology.

\sectionline

\subsection{Text}

\bvg\bva \alst{W}ola, \alst{w}iht, taz tú \alst{w}ęist, \hld\ taz tu \alst{w}iht hęizist, &
taz tú ne węist noch ne \alst{ch}anst \hld\ cheden \alst{ch}nos-pinci.\eva

\bvb Lo, wight, thou knowest that thou art called a wight, \\
that thou neither knowest nor canst say [the word] ‘Knosping’.\evb\evg

\sectionline
